% Preamble for "Lean for Working Algebraists" -- LaTeX manuscript,
% generated from lean_book/*.md by build/build_latex.py. Plain, portable
% LaTeX (no Springer class file yet -- that is a separate later decision;
% see the Phase 5 CHANGELOG entry).
\documentclass[11pt]{book}

\usepackage{fontspec}
\setmonofont{Consolas}
\setmainfont{Palatino Linotype}

\usepackage{amsmath}
\usepackage{amssymb}
\usepackage{amsthm}
\usepackage{mathtools}

\usepackage{tikz-cd}
\usepackage{graphicx}
\usepackage{longtable}
\usepackage{booktabs}
\usepackage{array}
\usepackage{calc}
\usepackage{caption}

\usepackage{listings}
\usepackage{xcolor}
% `listings` style definitions for Lean and Python code, matching each
% other in font/frame/color so the two read as one consistent system.
% Lean's keyword/tactic lists and comment syntax match Lean 4's actual
% grammar, adapted to `listings`' simpler (non-context-sensitive) model.
\lstdefinelanguage{lean}{
  morekeywords={theorem,lemma,def,abbrev,structure,inductive,class,instance,
    namespace,end,section,import,open,variable,variables,universe,universes,
    example,axiom,constant,by,where,with,match,fun,let,have,show,from,do,
    deriving,mutual,partial,private,protected,noncomputable,unsafe,extends,
    if,then,else,calc,set_option,attribute,macro,syntax,elab,notation,infix,
    infixl,infixr,prefix,postfix,this,Type,Prop,Sort},
  morekeywords=[2]{rw,rewrite,exact,apply,intro,intros,simp,simpa,simp_all,
    dsimp,cases,rcases,obtain,induction,constructor,refine,decide,omega,
    norm_num,ring,ring_nf,noncomm_ring,abel,field_simp,linarith,nlinarith,
    positivity,tauto,aesop,use,left,right,funext,ext,congr,unfold,show,
    sorry,rfl,trivial,assumption,contradiction,specialize,subst,change,
    conv,conv_lhs,conv_rhs,by_contra,exact?,apply?,gcongr},
  sensitive=true,
  morecomment=[l]{--},
  morecomment=[s]{/-}{-/},
  morestring=[b]{"},
}

\definecolor{codekeyword}{RGB}{20,60,150}
\definecolor{codetactic}{RGB}{110,40,140}
\definecolor{codecomment}{RGB}{60,110,60}
\definecolor{codestring}{RGB}{150,40,40}
\definecolor{codebg}{RGB}{247,247,247}

\lstdefinestyle{lean}{
  language=lean,
  basicstyle=\ttfamily\small,
  keywordstyle=\bfseries\color{codekeyword},
  keywordstyle=[2]\color{codetactic},
  commentstyle=\itshape\color{codecomment},
  stringstyle=\color{codestring},
  backgroundcolor=\color{codebg},
  showstringspaces=false,
  breaklines=true,
  breakautoindent=true,
  columns=flexible,
  frame=single,
  framerule=0.3pt,
  rulecolor=\color{black!30},
  numbers=none,
  tabsize=2,
}

\lstdefinestyle{python}{
  language=Python,
  basicstyle=\ttfamily\small,
  keywordstyle=\bfseries\color{codekeyword},
  commentstyle=\itshape\color{codecomment},
  stringstyle=\color{codestring},
  backgroundcolor=\color{codebg},
  showstringspaces=false,
  breaklines=true,
  breakautoindent=true,
  columns=flexible,
  frame=single,
  framerule=0.3pt,
  rulecolor=\color{black!30},
  numbers=none,
  tabsize=4,
}

% `style=lean` as the document-wide default (every fence in this book is
% either `lean` or `python`; build_latex.py rewrites Pandoc's emitted
% `[language=python]` to `[language=Python,style=python]` so Python fences
% pick up listings' built-in Python rules and this file's matching style,
% since `listings`' language names are case-sensitive and Pandoc emits the
% Markdown fence's info string verbatim, lowercase).
\lstset{style=lean}

% Pandoc's --listings output wraps every inline code span in
% \passthrough{\lstinline!...!}; this macro is normally supplied by
% Pandoc's own --standalone template, which this fragment-based pipeline
% does not use.
\providecommand{\passthrough}[1]{#1}
\providecommand{\tightlist}{\setlength{\itemsep}{0pt}\setlength{\parskip}{0pt}}
% Pandoc's longtable output computes column widths as e.g.
% "(\linewidth - 2\tabcolsep) * \real{0.5000}"; \real is normally supplied
% by Pandoc's own --standalone template (which this fragment-based
% pipeline does not use) as a trivial pass-through of its numeric argument.
\providecommand{\real}[1]{#1}
% Pandoc wraps every image in \pandocbounded{...} to constrain it to the
% text width/height; normally defined by Pandoc's --standalone template.
\providecommand{\pandocbounded}[1]{#1}

% Unicode symbols used verbatim inside Lean/Python `lstlisting` blocks.
% listings itself does not know these; newunicodechar maps each raw
% character, at the XeTeX input level, to the equivalent LaTeX command,
% independent of any particular font's glyph coverage.
\usepackage{newunicodechar}
\newunicodechar{§}{\S}
\newunicodechar{¬}{\ensuremath{\neg}}
\newunicodechar{·}{\ensuremath{\cdot}}
\newunicodechar{¹}{\textsuperscript{1}}
\newunicodechar{×}{\ensuremath{\times}}
\newunicodechar{Π}{\ensuremath{\Pi}}
\newunicodechar{Σ}{\ensuremath{\Sigma}}
\newunicodechar{α}{\ensuremath{\alpha}}
\newunicodechar{β}{\ensuremath{\beta}}
\newunicodechar{λ}{\ensuremath{\lambda}}
\newunicodechar{•}{\ensuremath{\bullet}}
\newunicodechar{⁻}{\ensuremath{{}^-}}
\newunicodechar{ₗ}{\ensuremath{{}_l}}
\newunicodechar{ⱼ}{\ensuremath{{}_j}}
\newunicodechar{∀}{\ensuremath{\forall}}
\newunicodechar{∃}{\ensuremath{\exists}}
\newunicodechar{∈}{\ensuremath{\in}}
\newunicodechar{∘}{\ensuremath{\circ}}
\newunicodechar{∣}{\ensuremath{\mid}}
\newunicodechar{∧}{\ensuremath{\wedge}}
\newunicodechar{∨}{\ensuremath{\vee}}
\newunicodechar{≃}{\ensuremath{\simeq}}
\newunicodechar{≠}{\ensuremath{\neq}}
\newunicodechar{≡}{\ensuremath{\equiv}}
\newunicodechar{≥}{\ensuremath{\geq}}
\newunicodechar{⊢}{\ensuremath{\vdash}}
\newunicodechar{⟨}{\ensuremath{\langle}}
\newunicodechar{⟩}{\ensuremath{\rangle}}
\newunicodechar{⟶}{\ensuremath{\longrightarrow}}
\newunicodechar{→}{\ensuremath{\rightarrow}}
\newunicodechar{←}{\ensuremath{\leftarrow}}
\newunicodechar{ℤ}{\ensuremath{\mathbb{Z}}}
\newunicodechar{—}{\textemdash}

% Custom environments for the book's recurring boxes.
\theoremstyle{remark}
\newtheorem*{mathreading}{Mathematical reading}
\newtheorem*{progcorner}{Programmer's corner (Python)}
\newtheorem{pblproject}{Checkpoint Project}[chapter]

\usepackage[backend=biber,style=numeric,sorting=nyt]{biblatex}
\addbibresource{references.bib}

\usepackage[colorlinks=true,linkcolor=blue,citecolor=blue,urlcolor=blue]{hyperref}
\usepackage[capitalize]{cleveref}
